\section{Discussion and open problems}\label{sec:discussion}

\paragraph{The dynamical picture.} Theorems~\ref{thm:A} and~\ref{thm:B} identify the upper Abel-averaged transport exponents of every supercritical type~I analytic operator, for every moment, at every phase of the full-measure set $G^v_\alpha$: all of them equal $\max\{0,1-\ell_*/\beta(\alpha)\}$ (Corollary~\ref{cor:allmoments}). For the almost Mathieu operator, Damanik and Tcheremchantsev proved that the lower transport exponents vanish for every $\lambda>1$ and every frequency~\cite{DT2007}, so the fastest lower and upper exponents differ in the window $1<\lambda<e^{\beta}$. The proofs suggest a picture of strong intermittency. Near the resonant times $T\approx|q_n\alpha-p_n|^{-1}$ of strongly Liouville scales, the part of the packet at energy $E$ makes a burst, moving with speed about $e^{-L(E)q_n}$ and reaching distance $T^{1-L(E)/\beta_n}$; a fixed fraction of the packet takes part (Theorem~\ref{thm:B}(1)). The exponents record the largest of these bursts. We prove nothing about the proportion of times spent in each regime, and nothing about the lower exponents of potentials other than the almost Mathieu operator.

\paragraph{The open hypothesis: weak separation beyond type~I.} Theorem~\ref{thm:D} reduces the law $\max\{0,1-\ell_*/\beta(s\alpha)\}$ for every real-analytic potential with positive Lyapunov exponent on its spectrum to one hypothesis, WXS$(w,s\alpha)$ for its primitive part $w$, and the hypothesis is needed only when $w$ has acceleration at least two at some point of the spectrum. The open problem is therefore exactly this: \emph{does WXS hold for primitive real-analytic potentials with acceleration at least two?} WXS is a horizontal condition. It asks that, at each fixed energy near the spectrum, no two zeros of the rational discriminant in the phase variable, real or complex, cluster at an exponential rate near the drifting phase $q_n\theta$. It is not a condition on vertical gaps between band centres at a fixed phase: by Remark~\ref{rem:bandproduct}, vertical gaps enter the horizontal loss only through a product over all bands. Whether exponentially small vertical gaps recur, at a positive exponential rate, along the convergents of a fixed irrational frequency, and whether they then force horizontal clusters, is part of the open question. For the two-harmonic family $2\lambda\cos2\pi x+2\mu\cos4\pi x$, Proposition~\ref{prop:EDelta} reduces WXS to the phase-free estimate (E-$\Delta$) on a single polynomial of degree at most $2q$, which we also leave open. A proof of WXS, or a fixed frequency at which the law fails together with the corrected arithmetic invariant, would complete the picture for real-analytic potentials. The law could also hold where WXS fails; our methods do not decide this.

\paragraph{Exceptional phases.} The upper bounds are proved for $\theta\in G^v_\alpha$, the phases whose local phase $q_n\theta$ stays polynomially far from the turning points of the band centres. At other phases a turning point near the origin changes the geometry of the discriminant, and the window estimate must be modified. Whether the exponent can be larger at such phases is open.

\paragraph{Lower local dimensions.} Theorem~\ref{thm:B}(3) determines the upper local dimension of $\mu_\theta$. The lower local dimension and the Hausdorff dimension of spectral measures are not addressed here.

\paragraph{Critical and mixed spectrum.} The law concerns potentials whose Lyapunov exponent is positive on the whole spectrum. When the Lyapunov exponent vanishes somewhere on the spectrum and $\beta(\alpha)>0$, transport is quasiballistic by the theorem of Jitomirskaya and Zhang~\cite[Theorem~3]{JZ2022}. At critical energies with $\beta(\alpha)=0$ our methods give nothing: the resonant scales disappear, and the speed bound $e^{-L(E)q}$ becomes vacuous.

\paragraph{Physical meaning.} Laboratory realizations of the Aubry--Andr\'e model use frequencies such as the golden mean, for which $\beta=0$; Corollary~\ref{cor:tI-beta0} confirms the absence of power-law transport for every analytic potential whose Lyapunov exponent is positive on the spectrum. The theorems show that this absence is arithmetic, not a consequence of positive Lyapunov exponent alone, and that it does not depend on the details of the potential. The proofs indicate that a frequency with one exceptionally large partial quotient after a denominator $q$ produces a transient burst over the time window $e^{L(E)q}\lesssim T\lesssim|q\alpha-p|^{-1}$, energy by energy. The exponents themselves are infinite-time quantities of an infinite system.

\subsection*{How this work was done}

The authors developed the results and proofs with assistance from AI models from more than one developer. They take responsibility for every claim. The authors chose the question, directed the work and decided which lines to pursue; the models did most of the exploration, derivation and drafting.

The proof was developed over several long working sessions, each building on the previous ones and recording its proofs, failed routes and self-checks. Each step was then checked by a separate session that had not seen the material, with instructions to refute it and to read the primary sources it relies on. The repairs found in these checks are incorporated here. Priority and literature searches were AI-assisted. No computation is used in any proof.
